\documentclass[11pt]{article}

\usepackage[T1]{fontenc}
\usepackage[margin=1.1in]{geometry}
\usepackage{amsmath}
\usepackage{amssymb}
\usepackage{booktabs}
\usepackage{array}
\usepackage{xcolor}
\usepackage{caption}
\usepackage{microtype}
\usepackage[colorlinks=true,linkcolor=blue!50!black,citecolor=blue!50!black,urlcolor=blue!50!black]{hyperref}

\hypersetup{
  pdftitle={The Detection Floor: What a Null Result Was Measured With},
  pdfauthor={Allen Hsieh}
}

\captionsetup{font=small, labelfont=bf}

% Float packing: let tables share pages with more text instead of drifting.
\renewcommand{\topfraction}{0.9}
\renewcommand{\bottomfraction}{0.6}
\renewcommand{\textfraction}{0.07}
\renewcommand{\floatpagefraction}{0.8}
\setcounter{totalnumber}{4}
\setcounter{topnumber}{3}

\newcommand{\code}[1]{\texttt{#1}}
\newcommand{\usff}{\code{us54}}

\title{The Detection Floor:\\ What a Null Result Was Measured With}
\author{Allen Hsieh\\[2pt] {\small FishAI project}}
\date{August 2026}

\begin{document}

\maketitle

\begin{abstract}
\noindent
A research codebase that publishes null results owes its readers a characterisation of the
instrument those nulls were measured with, because an unresolved effect and an absent effect look
identical in a table. The FishAI project published several nulls for the \usff{} dialect of
Literature (Canadian Fish) before anyone measured what its one instrument --- $K$ duplicate deal
pairs of arm $X$ against arm $Y$, reported as the mean paired set-difference $\pm 1.96\,\mathrm{SE}$
--- could resolve. This paper reports that characterisation and applies it to the project's own
back catalogue. The per-pair standard deviation is $3.15$--$3.44$ sets for arms that differ at every
decision ($3.18$--$3.44$ on an independent replication), giving an 80\%-power minimum detectable
effect running from $0.76$ sets at $N=150$ through $0.47$ at $N=400$, $0.33$ at $N=800$ and $0.14$ at
$N=4300$. Validated against a planted edge of continuously variable known size, the normal-theory
power model held in all twelve (effect, $N$) cells of both runs to within about one binomial standard
error, and the byte-exact control --- which had to be \emph{changed} before it tested anything, an
earlier form short-circuiting past the very wrapper whose inertness was the claim --- prints
$0.0000 \pm 0.0000$. Two parts did not replicate and are reported unsoftened: the empirical floor at
$N=400$ moved from $84.7\%$ detection to $75.0\%$ on fresh banks (pooled $81.7\%$), so the floor is
only bounded to $0.36$--$0.50$; and nominal-95\% interval coverage at $N=400$ went $6.1\%$
($11/180$) to $10.8\%$ ($13/120$), pooling to $8.0\%$ $[5.2, 11.7]$, which excludes the nominal
$5\%$ --- while $N=800$ is nominal in both runs, a pattern that does not scale as a real coverage
defect should, so the cause is unknown. We name the trap the first run fell into: the minimum
detectable effect depends on \emph{that cell's own} standard deviation, which is not constant
($3.645$ for one cell; $4.57$--$4.95$ implied by a foreign-engine harness), and a cell whose two arms
mostly agree can resolve $+0.2475 \pm 0.1024$ at $n=800$ while appearing to sit under a generic
$0.33$ floor. Applying the floor reclassifies two of this project's published nulls as
non-measurements. The recommendation, grounded only in this codebase's experience, is to publish the
floor beside the null.
\end{abstract}

\section{Introduction}
\label{sec:intro}

There are two reasons a measured effect can be reported as zero. The effect can be absent, or the
instrument can be unable to see it. In a results table these look the same: a small number, an
interval containing zero, and a row that reads as settled. The distinction is not pedantic. The
first licenses abandoning a mechanism; the second licenses nothing at all except a bigger
experiment. Altman and Bland put the point in five words that have not been improved on --- absence
of evidence is not evidence of absence \cite{altman}.

This paper is about a codebase that made the mistake and then measured its way out of it. The FishAI
project builds an engine, a bot laboratory and a results site for six-player Literature (Canadian
Fish) in its 54-card \usff{} dialect \cite{v05}. Its research documents are unusual in one respect
that makes the failure especially visible: they publish negative results as first-class findings,
in the body text, with the same typography as the positive ones \cite{concession, containedbook,
inertaxis}. Four of the project's six papers at the time of writing carry a negative headline.

Every one of those headlines --- positive and negative alike --- comes out of a single instrument:
$K$ duplicate deal pairs of arm $X$ against arm $Y$ under \usff{}, reported as the mean paired
set-difference $\pm 1.96\,\mathrm{SE}$. Until the work reported here, nobody had characterised what
that instrument can and cannot see. Every published null was therefore ambiguous between ``no
effect'' and ``no resolution'', and the project had no principled way to tell its readers which it
had found. Worse, the ambiguity is asymmetric in exactly the direction that flatters a project that
prides itself on negative results: an underpowered design manufactures publishable-looking nulls for
free.

The remedy is a scratch harness, \code{scripts/probe-floor.mjs} \cite{probefloor}, and an
independent agent's re-run of the whole thing on fresh seed banks. Together they answer four
questions --- what is the noise level, what is the smallest planted effect the design actually
detects, how often does a true null print an interval excluding zero, and which committed numbers
sit below their own floor. The results are recorded in \cite{concession}~\S8a, which this paper
reports.

Three of the four parts replicated. Two did not, and both cut against the first run's conclusions.
This paper reports the failures at the same size as the successes, because a floor paper that
softened its own non-replications would be a self-refuting document.

The paper is organised as follows. Section~\ref{sec:instrument} states the instrument precisely.
Section~\ref{sec:noise} gives the noise level and the resulting power curve.
Section~\ref{sec:planted} describes the planted-edge validation, and
Section~\ref{sec:control} --- the methodological heart of the paper --- describes why the byte-exact
control had to be changed before it was worth anything. Section~\ref{sec:nonrep} reports the two
parts that did not replicate. Section~\ref{sec:trap} names the substitution trap that the first run
fell into. Section~\ref{sec:catalogue} applies the floor to the project's committed results and
reclassifies two of its published nulls. Section~\ref{sec:limits} states what none of this
establishes, and Section~\ref{sec:rec} makes the one recommendation the evidence supports.

\section{The instrument}
\label{sec:instrument}

\subsection{What a datum is}

Literature is played by six players in two teams of three, seats $0,2,4$ against $1,3,5$
\cite{v05}. Under \usff{} the 54-card deck is partitioned into nine books of six cards; the game
ends the moment a team is awarded five books. A game therefore terminates in a clinch, and the
outcome variable used throughout the project's concession research is not the win indicator but the
\emph{net sets}: the acting team's score minus the opposing team's.

The instrument is duplicate pairing, borrowed from duplicate bridge and used so that the deal is
never a confound \cite{botlab}. Every seed is played twice, once with arm $X$ on team $0$ and once
with $X$ on team $1$, both arms otherwise identical. The pair's datum is
\begin{equation}
\label{eq:datum}
d_i \;=\; \bigl(X - Y \mid X \text{ on team } 0\bigr) \;+\; \bigl(X - Y \mid X \text{ on team } 1\bigr),
\end{equation}
on the same seed, and the reported statistic over $K$ pairs is
$\bar{d} \pm 1.96 \cdot \mathrm{SD}(d)/\sqrt{K}$ \cite{probefloor}. The constant $1.96$ is a house
constant, fixed so that every interval in the repository is comparable to every other.

Two properties of this design bear on everything below. First, a \emph{mirror match} --- the same
arm against itself --- returns exactly zero for every pair by construction, since the two
orientations cancel term by term. That is a useful correctness check and a useless noise estimate:
it tells you nothing about the variance of a comparison between arms that actually differ. Second,
the pairing removes the deal but not the play. Two arms that differ at any decision produce
different game trees on the same deal, and the residual variance of $d_i$ is what the floor is made
of.

\subsection{The four questions}

\code{probe-floor.mjs} \cite{probefloor} answers four questions, in four stages that can be run
independently:

\begin{enumerate}
  \item \textbf{\code{sd}} --- the per-pair standard deviation under the conditions the published
  tables actually used, and the analytic minimum detectable effect (MDE) curve computed off it. The
  default is $1{,}200$ pairs per condition.
  \item \textbf{\code{calib}} and \textbf{\code{detect}} --- the \emph{empirical} floor. An edge of
  continuously variable true size is planted, measured at large $N$ ($2{,}000$ pairs per grid
  point), and then the fraction of independent banks whose nominal 95\% interval excludes zero is
  counted at $N=400$ and $N=800$. The default is 36 banks of 800 pairs per planted size.
  \item \textbf{\code{fpr}} --- the false-positive rate, on true nulls that nevertheless move the
  game.
  \item \textbf{\code{published}} --- the floor applied to the numbers already committed.
\end{enumerate}

Every seed bank the probe uses is a fresh prefix under \code{floor-}, disjoint from all twelve
banks in use elsewhere in the repository. The single exception is deliberate and is described in
Section~\ref{sec:control}.

\subsection{The power arithmetic}

The MDE is the standard two-sided normal-theory quantity for a paired mean at 80\% power:
\begin{equation}
\label{eq:mde}
\mathrm{MDE}(N) \;=\; \frac{\bigl(z_{0.975} + z_{0.80}\bigr)\,\mathrm{SD}}{\sqrt{N}}
\;=\; \frac{2.8016 \cdot \mathrm{SD}}{\sqrt{N}} ,
\end{equation}
with $z_{0.975} = 1.96$ the house constant and $z_{0.80} = 0.8416$ the one-sided power point
\cite{probefloor, cohen}. Nothing about equation~\eqref{eq:mde} is novel; the contribution here is
that the $\mathrm{SD}$ in it was, in this repository, never measured before it was needed, and that
when it was measured it turned out not to be a constant (Section~\ref{sec:trap}).

\section{The noise level, and the curve that follows}
\label{sec:noise}

\subsection{Per-pair standard deviation}

Table~\ref{tab:sd} gives the per-pair SD for each condition the probe measures, as recorded in the
harness's own \code{MEASURED\_SD} constant --- the values its \code{sd} stage returned on
2026-08-31 at $1{,}200$ pairs per condition, checked in so the \code{published} stage can be re-run
without paying for the sweep again \cite{probefloor}.

\begin{table}[ht]
\centering
\caption{Per-pair SD of the duplicate-pair difference, by condition, $1{,}200$ pairs each
\cite{probefloor}. The first eight rows are arms that differ at every decision; the last two are
arms that agree almost everywhere, and are included here because Section~\ref{sec:trap} turns on
the difference.}
\label{tab:sd}
\small
\begin{tabular}{lr}
\toprule
condition & per-pair SD (sets) \\
\midrule
defuse, balanced (\cite{concession}~\S3.1) & 3.321 \\
defuse, ghost (\S3.1, narrowest) & 3.237 \\
defuse, archivist (\S3.1, widest) & 3.428 \\
conceal vs.\ defusing opponent (\S5a.3) & 3.428 \\
conceal vs.\ non-defusing opponent (\S5a.3) & 3.363 \\
lopsided: balanced vs.\ turtle & 3.438 \\
gauntlet difference-of-differences vs.\ balanced (\S3.3) & 3.364 \\
information-free null, $\lambda = 0.25$ (\S2.1) & 3.150 \\
\midrule
planted edge, $\varepsilon = 0.02$ & 2.549 \\
symmetric null, $\varepsilon = 0.02$ & 2.959 \\
\bottomrule
\end{tabular}
\end{table}

The headline is the range over the eight full-divergence conditions: \textbf{$3.15$ to $3.44$ sets
per pair}. The independent replication on fresh banks returned $3.18$--$3.44$ \cite{concession}.
This is the most stable result in the whole exercise, and it is worth pausing on why it should be
stable: it is a property of the game's outcome distribution under duplicate pairing, not of any
particular mechanism, and it varies remarkably little across conditions as different as a one-knob
ablation and a lopsided style mismatch.

\subsection{The curve}

Substituting into equation~\eqref{eq:mde} gives Table~\ref{tab:mde}, the operative design tool of
this paper.

\begin{table}[ht]
\centering
\caption{80\%-power minimum detectable effect, in sets per duplicate pair, at the measured per-pair
SD \cite{concession}.}
\label{tab:mde}
\begin{tabular}{lrrrrrrr}
\toprule
$N$ (pairs) & 150 & 300 & 400 & 600 & 800 & 2000 & 4300 \\
\midrule
MDE (sets/pair) & 0.76 & 0.54 & 0.47 & 0.38 & 0.33 & 0.21 & 0.14 \\
\bottomrule
\end{tabular}
\end{table}

Read it as a constraint rather than a licence. At $N=400$ --- the size at which several of this
project's cells were run --- an effect of a third of a set is more likely than not to be missed. At
$N=150$, the size of the cross-engine cells in \cite{crossplay}~\S2, the design cannot resolve
anything under three quarters of a set.

\section{Validating the model: an edge of known size}
\label{sec:planted}

A power curve computed from a measured SD is a prediction, not a measurement. It assumes normality
of $\bar{d}$, correct SD estimation, and independence across pairs. The probe therefore checks the
prediction against an edge whose true size is known because it was put there.

\subsection{The planted edge}

The handicapped arm plays the shipped policy except that, on an ordinary ask, it replaces the chosen
ask with a uniformly chosen legal alternative with probability $\varepsilon$ \cite{probefloor}. Two
properties make this a usable ruler:

\begin{itemize}
  \item at $\varepsilon = 0$ no random draw is ever taken --- \code{u01} returns a value in $[0,1)$
  and every such value fails the $\geq \varepsilon$ test --- so the arm is byte-identical to the
  control and the paired difference must print exactly $0.0000 \pm 0.0000$;
  \item the true edge grows monotonically with $\varepsilon$, because replacing an argmax by a
  uniform draw over the legal ask list can only lose value in expectation.
\end{itemize}

The draw is a deterministic hash of (stream, seed, step, seat), so a bank replays identically and
the two orientations of a pair remain the same deal --- which is what keeps
equation~\eqref{eq:datum} meaningful under the handicap. The calibration grid is
$\varepsilon \in \{0.005, 0.01, 0.02, 0.03, 0.04, 0.06, 0.08\}$, measured at $2{,}000$ pairs per
point so that the true edge at each $\varepsilon$ is itself known far more precisely than the
$N=400$ and $N=800$ banks that are then asked to detect it.

\subsection{Two nulls that move the game}
\label{sec:nulls}

A mirror match is exactly zero by construction (Section~\ref{sec:instrument}), so it cannot serve as
the null against which false positives are counted. The probe uses two true nulls that
\emph{do} disturb play \cite{probefloor}:

\begin{itemize}
  \item \textbf{\code{symnull}} --- both arms carry the \emph{same} handicap magnitude
  $\varepsilon$, on two different hash streams. They differ at roughly $2\varepsilon$ of all ask
  decisions, yet the expected edge is exactly zero by the exchangeability of the two streams. This
  is the cleanest null available: the game moves, and the truth is nonetheless known.
  \item \textbf{\code{abnull}} --- the information-free control of the project's earlier probe
  \code{probe-ab.mjs}, reused as specified: the danger penalty's magnitude with the threat estimate
  $D$ replaced by a hash of (seat, log length) over the same $0$--$3$ prey range. The prototype
  ranker is copied faithfully, so that appetite $\lambda = 0$ reproduces \code{decide} at
  \code{defuse: 0} byte-for-byte. That is a weaker statement than reproducing the shipped decision
  function, and the difference matters: the shipped roster carries \code{defuse: 1}, so what this
  control certifies is fidelity to the \emph{unequipped} baseline, not to the policy the project
  ships \cite{probefloor}.
\end{itemize}

In the \code{fpr} stage \code{symnull} is run at two magnitudes, $\varepsilon = 0.02$ and
$\varepsilon = 0.08$, and \code{abnull} at the single appetite $\lambda = 0.25$ that
\code{probe-ab.mjs} used \cite{probefloor}. All three arms are run at banks of 800 pairs each split
into its two disjoint halves of 400, so that each bank yields one estimate at $N=800$ and two at
$N=400$.

\subsection{The power model holds}

This is the solid finding of the exercise. Across twelve (planted-effect, $N$) cells --- six planted
sizes at two sample sizes --- the observed detection rate sat within about one binomial standard
error of the normal-theory prediction in \emph{both} runs \cite{concession}. Equation~\eqref{eq:mde}
can therefore be trusted as a design tool: an experimenter in this repository who wants 80\% power
against a half-set effect can read off the required $N$ and expect to get what the arithmetic
promises.

What the model does \emph{not} deliver is a tight empirical estimate of where the floor sits, which
is Section~\ref{sec:nonrep}'s subject and a different question. The model being calibrated in
aggregate is compatible with the single grid point nearest the floor being estimated badly, and that
is what happened.

\section{The control that had to be changed before it meant anything}
\label{sec:control}

The project's probe discipline requires a byte-exact control: an arm which, with the mechanism
switched off, reproduces the shipped policy exactly, so that the harness prints
$0.0000 \pm 0.0000$ and any nonzero reading elsewhere is attributable to the mechanism rather than to
the scaffolding. Every published concession result in \cite{concession} carries one.

The floor probe's control is the planted-edge arm at $\varepsilon = 0$. An earlier version of that
arm opened with a short-circuit:

\begin{quote}
\begin{verbatim}
if (!(eps > 0)) return plainPolicy(style)
\end{verbatim}
\end{quote}

\noindent
which is the natural thing to write and is worthless as a control. With the short-circuit in place
the control returns before the wrapper does anything, so the comparison it certifies is
\code{plainPolicy(style)} against the shipped policy --- a statement true by inspection, since
\code{plainPolicy} is a one-line delegation to \code{decide}. The wrapper whose inertness at
$\varepsilon = 0$ is the actual claim is never entered. The draw is never taken, the legality lookup
never runs, the substitution branch is never reached, and a bug in any of the three would sail
through a control printing a perfect zero \cite{probefloor}.

The fix is to delete the short-circuit and let the arm fall through. The arm then takes the draw,
computes the legal ask list, evaluates the substitution branch's guard, declines because
$u01 \geq 0$ always holds, and \emph{still} prints $0.0000 \pm 0.0000$ over 200 pairs. That is the
statement the control discipline was asking for. The harness's own comment states the point
plainly, and the file preserves the superseded version in prose so the reasoning survives the diff.

Three observations generalise beyond this file.

\textbf{A control must exercise the code whose inertness it certifies.} The failure mode is not that
the short-circuited control was wrong --- its output was correct --- but that it was
\emph{vacuous}: it verified a proposition nobody doubted while appearing to verify one that mattered.
A green control is not informative unless one can say which lines it ran.

\textbf{The fast path and the null path are the same path here.} Optimising a null case by returning
early is normally harmless and often good practice. It is specifically harmful when the null case is
the control, because the optimisation removes exactly the coverage the control exists to provide.
Any codebase whose experimental controls are ``the mechanism at appetite zero'' has this hazard
structurally.

\textbf{The fidelity check is a separate obligation.} A control shows the scaffolding is inert; it
does not show the scaffolding measures the same thing the published harness measured. The probe
therefore adds a fourth control job that deliberately breaks its own fresh-bank rule and replays
\code{holdoutA}, the bank on which \cite{concession}~\S3.1 published
$+1.4250 \pm 0.3179$ over 400 pairs for the shipped defusal mechanism. Its purpose is to reproduce a
published number rather than to measure a new one, and that is the only reason a probe should ever
reuse a bank \cite{probefloor}.

\section{The two parts that did not replicate}
\label{sec:nonrep}

An independent agent re-ran the entire probe on fresh seed banks. Two of the four parts came back
different, and both differences remove a claim the first run had made.

\subsection{The empirical floor at \texorpdfstring{$N=400$}{N=400}}
\label{sec:floor400}

The first run found that a planted true effect of $0.393$ sets per pair was detected --- its nominal
95\% interval excluding zero --- in $84.7\%$ of independent banks at $N=400$. That is comfortably
above the $80\%$ target, and the first run drew the conclusion that the empirical floor
\emph{matched the analytic MDE to 1.5\%}: a satisfying result, and the kind of agreement between
prediction and measurement that invites trust in both.

The replication detected the same planted effect in $75.0\%$ of banks. Pooled over both runs, the
detection rate at that grid point is $81.7\%$. Table~\ref{tab:nonrep} sets the two runs side by
side.

\begin{table}[ht]
\centering
\caption{The two non-replications \cite{concession}. Every figure is as printed in the source; the
interval construction used for the coverage rows is not recorded alongside them.}
\label{tab:nonrep}
\small
\begin{tabular}{lccc}
\toprule
quantity & first run & replication & pooled \\
\midrule
detection of a $0.393$ effect at $N=400$ & 84.7\% & 75.0\% & 81.7\% \\
empirical floor at $N=800$ & 0.256 & 0.275 & replicates \\
coverage: intervals excluding zero, $N=400$ & 6.1\% (11/180) & 10.8\% (13/120) & \textbf{8.0\% (24/300)} \\
\quad 95\% CI on that rate & --- & $[5.9, 17.8]$ & $[5.2, 11.7]$ \\
coverage at $N=800$ & nominal & nominal & 4.0\% \\
\bottomrule
\end{tabular}
\end{table}

The consequence is that \textbf{the $N=400$ floor is bounded only to the interval
$0.36$--$0.50$ sets per pair, and this design cannot pin it tighter}. The first run's 1.5\%
agreement was an artifact of where the sweep grid happened to land: with a grid of planted sizes and
a detection rate estimated from a few dozen banks per point, the identity of the smallest point
clearing 80\% is a coarse and unstable statistic, and reading a two-significant-figure agreement off
it was over-reading. The analytic MDE at $N=400$, $0.47$, lies inside the bracket, which is
reassuring and is not the same as being confirmed.

The $N=800$ floor \emph{does} replicate: $0.256$ against $0.275$. The design is better behaved where
it has more data, which is what one would expect and is a reason to prefer the $N=800$ column of
Table~\ref{tab:mde} when a choice is available.

\subsection{Interval coverage at \texorpdfstring{$N=400$}{N=400}}
\label{sec:coverage}

This is the more uncomfortable of the two, because it removes the first run's strongest claim.

On true nulls that still move the game (Section~\ref{sec:nulls}), the first run counted $11$ of
$180$ nominal-95\% intervals excluding zero at $N=400$: $6.1\%$, which is nominal. On that basis the
first run concluded that \emph{the nominal 95\% intervals are correctly sized; no interval in this
repository needs widening}. That is a strong and useful claim, and it would have retired an entire
class of worry about the project's published error bars.

The replication counted $13$ of $120$: $10.8\%$, with a 95\% interval of $[5.9, 17.8]$ that excludes
$5\%$. Pooled over both runs the rate is $\mathbf{8.0\%}$ ($24/300$), with interval
$[\mathbf{5.2}, \mathbf{11.7}]$ --- which also excludes the nominal $5\%$. \textbf{The first run's
strongest claim fails on independent banks.}

At $N=800$ both runs are nominal, pooling to $4.0\%$.

We report the cause as \textbf{unknown}, and say why rather than leaving it at a shrug. A genuine
under-coverage defect --- a heavy-tailed pair distribution, a variance estimate biased low, a
residual dependence between the two orientations of a pair --- ought to scale with $N$ in a
recognisable way: worse at smaller $N$, tending to nominal as the central limit theorem takes hold.
What is observed is nominal at $N=800$ and anomalous at $N=400$, from banks that are literally the
halves of the $N=800$ banks. That is a very specific pattern and no mechanism proposed so far
predicts it. The two candidate readings are that $N=400$ is genuinely below the size at which the
normal approximation is safe for this distribution and $N=800$ is above it --- a sharp transition
over one doubling, which is a lot to ask --- or that the excess is a fluctuation and $24/300$ with an
interval whose lower end sits at $5.2\%$ is simply a near-miss dressed as a finding.

We decline to choose. What can be said without choosing is the practical instruction: \textbf{a
nominal 95\% interval computed at $N=400$ in this repository should not be treated as a 95\%
interval}, and the honest reading of a marginal $N=400$ result is that its error bar may be
optimistic by an unquantified amount. Whether this should be assumed benign is exactly the question
\cite{concession}~\S9 refuses to close.

\section{The substitution trap}
\label{sec:trap}

The first run of this probe fell into a trap worth naming, because it is the most transferable
mistake in the paper and because it is invisible if one reads Table~\ref{tab:mde} as a threshold
rather than as a formula.

\subsection{The SD is not a constant}

Equation~\eqref{eq:mde} contains a standard deviation. Table~\ref{tab:mde} was computed with one
particular value of it, from one particular family of conditions. The MDE of any given cell depends
on \emph{that cell's own} SD, and the SD is not constant across this repository:

\begin{itemize}
  \item The eight full-divergence conditions of Table~\ref{tab:sd} span $3.15$--$3.44$, which is
  narrow enough to invite the substitution.
  \item The Turtle cell of \cite{concession}~\S3.1 measures $\mathbf{3.645}$ --- outside that band,
  in the direction that makes its floor higher than the generic figure suggests.
  \item The cross-engine harness of \cite{crossplay}~\S3, which runs FishAI's policy inside a
  third-party engine against a foreign bot family, publishes half-widths of $0.517$, $0.555$,
  $0.533$ and $0.560$ at 300 duplicate pairs per opponent. Those imply a per-pair SD of
  $\mathbf{4.57}$ to $\mathbf{4.95}$: \textbf{38--49\% higher} than the \usff{} figure
  \cite{concession}.
\end{itemize}

Substituting one SD across all conditions understates the floor exactly where the cells are
smallest, which is exactly where the understatement does the most damage. The cross-engine cells are
the clearest case: they are both the noisiest and the shortest-run measurements in the project.

\subsection{The counter-example that keeps the reader honest}
\label{sec:counter}

The trap has a mirror image, and it matters more, because it is the error a reader makes with
Table~\ref{tab:mde} in hand and good intentions.

A head-to-head between two adaptive generations, \emph{v1.0 against v0.5} --- two members of the
same generational ladder \cite{v15} --- is recorded as $+0.2475 \pm 0.1024$ at $n=800$. Against the generic $0.33$ floor at $N=800$ that reads as ``below
the floor, not resolved'' --- and the reading is wrong. Its own interval excludes zero comfortably,
by a margin of more than two half-widths. The reason is structural: the generic floor was computed
from arms that differ at \emph{every} decision, and two adaptive arms that mostly agree produce a far
smaller per-pair SD than a one-knob ablation does, so their cell's own floor is far lower than the
generic one.

Table~\ref{tab:sd}'s last two rows are the committed evidence for that mechanism inside this
harness. The planted-edge arm at $\varepsilon = 0.02$ --- which differs from its control on about
2\% of ask decisions --- has a per-pair SD of $2.549$, and the symmetric null at the same magnitude
$2.959$, against $3.15$--$3.44$ for arms that differ everywhere. Agreement between arms buys
precision.

The rule that follows is short and is stated in \cite{concession} in bold for a reason:
\textbf{never read the floor table as a threshold to apply by eye}. The floor is a function with an
argument, and the argument is the cell's own noise.

\subsection{The counter-example is itself not traceable}
\label{sec:untraceable}

Honesty about a floor paper requires one more sentence about that counter-example, and it is not a
comfortable one. \textbf{The figure $+0.2475 \pm 0.1024$ cannot be found in any committed document.}
It appears exactly twice in the repository: in \cite{concession}~\S8a, where it does the work
described above, and in the probe's own list of published claims, where its provenance field reads
\code{'task brief'} and its description carries the parenthesis \emph{``NOT FOUND in any committed
doc''} \cite{probefloor}.

We report this rather than quietly substituting a traceable example, for three reasons. It is the
project's own rule that every number must be traceable to a committed source, and this one is not.
The probe's author noticed and labelled it, which is the behaviour the rule is meant to produce, and
suppressing the label would waste that. And the pedagogical point the number carries --- that two
mostly-agreeing arms resolve effects a generic floor would reject --- does not depend on it: the
mechanism is demonstrated inside the harness by Table~\ref{tab:sd}'s last two rows, on banks the
probe generated and can regenerate. The counter-example should be re-derived from a cell that exists
in a committed artifact, and until it is, it should be read as an illustration rather than as a
measurement.

\section{Applying the floor to the back catalogue}
\label{sec:catalogue}

The point of characterising an instrument is to go back and re-read what it produced.
Table~\ref{tab:catalogue} is \cite{concession}~\S8a.1: every headline result of the project's
concession research placed against its own floor.

\begin{table}[ht]
\centering
\caption{Committed results against their own detection floor \cite{concession}. ``Its floor'' is the
80\%-power MDE at that cell's $N$ and, where it could be measured, that cell's own SD.}
\label{tab:catalogue}
\small
\begin{tabular}{p{0.30\textwidth} r r r p{0.24\textwidth}}
\toprule
result & reported & $N$ & its floor & verdict \\
\midrule
\S5a.3 concealment vs.\ non-defusing opponents & $-0.1483 \pm 0.2571$ & 600 & ${\sim}0.38$ &
  \textbf{below floor --- not resolved} \\
\S6 the contained-pass aim, ``inert'' & $-0.045 \pm 0.130$ & --- & --- &
  \textbf{below floor --- a non-measurement, not a zero} \\
\cite{crossplay}~\S2, all four headline cells & 42.7--48.7\%, no intervals published & 150 &
  1.05--1.13 & \textbf{all four below floor} \\
\S3.1 defusal, shipped & $+1.43$ / $+1.58 \pm 0.32$ & 400 & ${\sim}0.47$ & resolved, $3\times$ the
  floor \\
\S8.1 the six triangle edges & $\pm 0.10$ half-widths & 4000 & ${\sim}0.15$ & five resolved; \emph{P}
  vs.\ \emph{C} ($+0.1457$, $+0.1742$) sits at the floor --- \textbf{nominal only} \\
\cite{bounded}~\S5a, v1.5 vs.\ v1.0 & $-0.1255 \pm 0.0590$ & see \S\ref{sec:discrep} & see note &
  resolved --- that cell's arms differ far less, so its own SD is much smaller \\
\bottomrule
\end{tabular}
\end{table}

\subsection{Two nulls become non-measurements}

The two reclassifications are the substantive output of this exercise, and both are corrections to
claims this project published.

\textbf{\cite{concession}~\S6's ``inert'' contained-pass aim.} The project replaced the aim and cost
model of its contained-book turn-pass --- the one place in the engine where ``concede to the least
dangerous seat'' is exactly right, because a contained-book pass is a guaranteed miss by
construction and so cannot suffer the danger/opportunity confound that sinks the same rule on
ordinary asks --- with the measured threat estimate, and got $-0.045 \pm 0.130$ over 400 duplicate
pairs. The reading is now \textbf{a statement about this instrument's resolution and not about the
game}: at the $N=400$ that cell was run at, the floor is $0.47$ sets per pair, so an effect of
$0.045$ is an order of magnitude under the only floor that applies to it --- and it stays a factor of
three under $0.14$, the smallest floor anywhere in Table~\ref{tab:mde}.

The reclassification is not a call to un-revert, and it is worth being exact about what it does and
does not touch, because the reason \S6 gives for the revert was never the interval.
\cite{concession}~\S6 states that the change was \textbf{reverted rather than shipped off by
default}, on that document's own \S6.2 argument: a knob that is listed, swept and reported while
being inert contaminates the measured vector, and one such knob is already one too many
\cite{inertaxis}. That reason stands untouched here. Two further things should be said in the
reverted change's defence. First, the diagnostic that accompanied the null gives an independent
mechanism for it: over $25{,}654$ ask decisions, more than one opponent still held cards $99.1\%$ of
the time, the contained pass fired at all on only $1.06\%$, and it fired \emph{and} the two aims
disagreed on $0.11\%$ \cite{concession}. A mechanism that gets to matter on one decision in a
thousand cannot move a set-difference, and that argument does not depend on the interval. Second,
the reverted aim was the soundest available form of the fix, built at the one place in the engine
where the danger/opportunity confound cannot arise. What the reclassification removes is only the
supporting reading --- the word ``inert'' --- and not the ground the revert was actually made on.

\textbf{\cite{concession}~\S5a.3's non-defuser cell.} Concealment measured $-0.1483 \pm 0.2571$ at
$N=600$ against opponents that do not defuse, and the document originally read that as concealment
being worth nothing against them. At that cell's own SD of $3.363$ the floor is about $0.38$ sets:
the measurement is not capable of distinguishing a real loss of about a third of a set from a small
gain. The honest statement, now in the source document, is that the cell is \textbf{not resolved}.
What survives is the comparative result measured at $4{,}000$ pairs per cell, where concealment-only
is the weakest of the four concession arms outright --- a claim whose floor is about $0.15$ and which
clears it.

\textbf{And a third that is not this project's own.} \cite{crossplay}~\S2's per-opponent ranking of
FishAI against a foreign bot family --- four headline cells at $42.7\%$ to $48.7\%$ win rate, at 150
games per cell, published without intervals --- sits under a floor of $1.05$--$1.13$ sets per pair on
that harness's own noise level. The ordering of those four opponents is not resolved. The
\emph{aggregate} claim that harness was built to make, that defusal replicates in a foreign engine at
$+1.55$ to $+1.75$ sets per pair over 300 pairs per opponent with four intervals excluding zero, is
untouched by this.

\subsection{Two discrepancies found while reproducing the table}
\label{sec:discrep}

Reproducing Table~\ref{tab:catalogue} against its underlying sources surfaced two bookkeeping
discrepancies. Neither changes a verdict; both are reported because a floor paper that quietly
normalised its own inputs would be exactly the kind of document it is arguing against.

\textbf{The bounded-memory row's $N$.} \cite{concession}~\S8a.1 records the v1.5-against-v1.0 cell at
$N=2700$. \cite{bounded}~\S5a, the document that published it, states $2{,}000$ duplicate pairs per
rung on the held-out bank, and the probe's own list of published claims carries $n = 2000$ for the
same row \cite{probefloor}. Two committed sources say $2{,}000$ and one says $2{,}700$; the balance
of evidence is that the table's $N$ is a transcription error. The row's verdict does not turn on it,
because that row is resolved for a reason about its SD rather than its $N$
(Section~\ref{sec:counter}) --- which is the only reason this is a footnote rather than a
correction.

\textbf{The untraceable counter-example}, Section~\ref{sec:untraceable}.

\subsection{A verdict rule worth stating separately}

The probe's \code{published} stage does not treat ``the published interval excluded zero'' as
sufficient for ``resolved''. It requires both that the interval excluded zero \emph{and} that the
reported effect is at or above the 80\%-power floor, and it prints a distinct verdict ---
\code{NOMINAL ONLY --- underpowered} --- for a cell that clears its own interval but not its floor
\cite{probefloor}.

The reasoning is standard and under-applied. An effect that clears its interval but sits below the
design's floor was found by a design that would have missed it more often than not; conditional on
being found at all, such an estimate is biased upward in magnitude, and if the true effect is small
its sign is not reliable either \cite{gelmancarlin}. Reporting it as ``significant'' without
reporting that its design had, say, 40\% power against it is how a literature accumulates effect
sizes that shrink on replication \cite{button, ioannidis}; the same argument is why the reform
literature asks for estimates and intervals to be read as continuous evidence rather than sorted
into significant and not \cite{amrhein}. Naming the underpowered-but-significant
case in the harness's own output, rather than in a caveat somewhere else, is the cheapest available
guard against it.

The rule has a cost, and this paper pays it on its own table rather than exempting itself. Applied
to \S8.1's six triangle edges, five clear the ${\sim}0.15$ floor at $N=4000$ by a wide margin, but
the weakest --- \emph{P} against \emph{C}, plain against conceal-only --- is $+0.1457 \pm 0.1030$ on
the first bank set and $+0.1742 \pm 0.1030$ on the replication \cite{concession}. Both intervals
exclude zero; neither effect is three times its floor, and the first sits fractionally
\emph{under} it. By this paper's own rule that edge is \textbf{nominal only, not resolved}. Nothing
about the triangle's refutation depends on it --- that argument is carried by wrong-signed legs many
times their floor --- but the ordering claim that conceal-only is the worst of the four arms rests on
this edge, and it rests on it at the floor.

\section{What this does not establish}
\label{sec:limits}

\textbf{The floor is a property of this instrument on this game.} $3.15$--$3.44$ sets per pair is
the residual noise of duplicate-paired net sets under \usff{} at this roster's level of play. It is
not a property of Literature, of duplicate pairing in general, or of any other outcome variable. The
cross-engine harness's $4.57$--$4.95$ is the nearest available demonstration that the number does not
transfer even to a different implementation of the same game.

\textbf{The floor applies to the paired set-difference and to nothing else.} Several of this
project's nulls are reported in other units. The contained-book paper's headline mirror table, for
instance, reports mean score on a $0.5$-centred scale with standard errors of $0.0000$--$0.0065$
\cite{containedbook}; that is a different statistic on a different scale, and Table~\ref{tab:mde}
says nothing about it. A floor characterisation is per-instrument, and this project has more than one
instrument.

\textbf{Nothing here rescues a null that is below its floor.} Reclassifying \S6 and \S5a.3 as
non-measurements does not make either mechanism valuable. It removes a claim; it does not install
the opposite claim. Both cells would need re-running at an $N$ chosen against the effect size
actually worth detecting, and neither has been.

\textbf{The empirical floor at $N=400$ is bracketed, not measured.} Section~\ref{sec:floor400}'s
$0.36$--$0.50$ is the honest width of what two runs of this design can say. Pinning it tighter needs
either many more banks per grid point or a finer grid, and the exercise would have to be re-costed
against what a tighter number would actually buy.

\textbf{The coverage anomaly is unresolved and may be nothing.} Section~\ref{sec:coverage} states
both readings and chooses neither. A reader who wants a single number should take $N=800$'s $4.0\%$,
use $N \geq 800$ for anything marginal, and treat $N=400$ intervals as provisional.

\textbf{The probes are not the lab.} \code{probe-floor.mjs} and its siblings are scratch harnesses:
not typechecked, emitting no committed artifact, and their numbers are not digest-bound
\cite{concession}. They exist so that every table can be regenerated, and that is a weaker guarantee
than the one the lab's committed results carry. The one mitigation is the fidelity check of
Section~\ref{sec:control}, which shows this harness reproduces a published number on a published
bank.

\textbf{Decision-level identity is a stronger tool where it applies.} Where the question is whether
two policies differ \emph{at all}, an exact comparison of their actions on identical observations and
seeds returns a certificate rather than an interval, and no floor arithmetic is needed
\cite{inertaxis}. The floor characterised here is the right instrument for ``how much is this worth''
and the wrong one for ``does this fire''. The project has both, and part of reading a null correctly
is noticing which question was asked.

\section{The recommendation}
\label{sec:rec}

One recommendation follows from this experience, and it is deliberately modest, because it is
grounded in one codebase, one game, and one instrument.

\textbf{Publish the floor beside the null.} A reported null should carry, in the same table row, the
smallest effect the design that produced it would have detected 80\% of the time --- computed with
that cell's own noise level, not a repository-wide constant. Nothing else in this paper is a
recommendation to anyone. This one is, for four reasons the work made concrete:

\begin{enumerate}
  \item \textbf{It is the cheapest possible remedy for the ambiguity.} The floor is one number,
  computed from a standard deviation the harness has already estimated in the course of printing its
  interval. It requires no additional games. Two of this project's nulls were misread for want of a
  number that was, in effect, already in memory.
  \item \textbf{It changes what the null licenses.} ``Measured at $-0.045 \pm 0.130$'' invites the
  reader to abandon the mechanism. ``Measured at $-0.045 \pm 0.130$, against a floor of $0.47$''
  invites the reader to note that nothing was learned and to ask what $N$ would be needed. Those are
  different actions, and only the second is supported.
  \item \textbf{It has to be per-cell, and stating it per-cell is what stops the substitution
  trap.} A repository-wide floor printed once in a methods section becomes a threshold applied by
  eye, which is precisely the error of Section~\ref{sec:trap}. Printed in the row, beside its own
  interval, the floor carries its argument with it.
  \item \textbf{It composes with the discipline of putting the correction where the claim is.} This
  project's practice is to correct claims adjacent to the claims themselves rather than in a
  footnote or a later document \cite{inertaxis, concession}. A floor in the row is that practice
  applied pre-emptively: the caveat cannot become detached from the number it qualifies, because it
  is part of the number's presentation.
\end{enumerate}

The negative form of the recommendation is worth stating too, because this paper is an instance of
it. \textbf{Do not publish a floor without replicating it.} The first run of this exercise produced
a clean, quotable, mutually confirming set of four results, and one independent re-run on fresh banks
removed two of them --- including the strongest. Had the floor been published after one run, this
project would have added a false reassurance about its own error bars to the pile of claims it was
trying to clean up. A characterisation of an instrument is itself a measurement, and it inherits
every obligation the measurements it is meant to qualify are held to.

\section{Conclusion}

An unresolved effect and an absent effect look identical in a table, and a project that publishes
nulls as findings is precisely the project that cannot afford the confusion. FishAI published
several before measuring what its instrument could resolve. Measured, the instrument's noise is
$3.15$--$3.44$ sets per duplicate pair, giving an 80\%-power floor of $0.47$ sets at $N=400$ and
$0.33$ at $N=800$; the normal-theory power model behind those figures was validated against a
planted edge of known size in all twelve cells of both runs, on the strength of a byte-exact control
that had first to be repaired, because the version that short-circuited at zero handicap proved only
that the unwrapped policy equalled the shipped policy and never entered the wrapper at issue. Two
parts did not replicate: the empirical floor at $N=400$ is bounded only to $0.36$--$0.50$, and
interval coverage at $N=400$ pools to $8.0\%$ $[5.2, 11.7]$ against a nominal $5\%$, for reasons
unknown and not assumed benign. The floor must always be that cell's own --- one cell measures
$3.645$, a foreign-engine harness $4.57$--$4.95$, and two mostly-agreeing arms resolve effects a
generic floor would reject. Applied to the back catalogue, the floor turns two published nulls into
non-measurements; the mechanisms the project ships clear their floors by a factor of three or more,
with one exception the same rule catches --- the weakest of \S8.1's six triangle edges, plain against
conceal-only, whose $+0.1457$ and $+0.1742$ sit at a floor of ${\sim}0.15$ and are nominal only.

The deliverable is not the number. It is the row the number belongs in: publish the floor beside the
null.

\begin{thebibliography}{99}
\small

\bibitem{concession}
FishAI project.
\emph{CONCESSION.md --- FishAI v2.0: the three-sided ask}.
Repository document; \S3.1 (defusal, shipped), \S5a.3 (concealment, and the non-defuser cell), \S6
(the contained-pass aim, reverted), \S8.1 (the triangle at 4{,}000 pairs), \textbf{\S8a} (the
detection floor, its replication, and the substitution trap), \S8a.1 (the floor applied to committed
numbers), \S9 (open residues), \S10 (reproducing the measurements).

\bibitem{probefloor}
FishAI project.
\emph{\code{scripts/probe-floor.mjs} --- the detection-floor harness}.
Source file; the four stages (\code{sd}, \code{calib}/\code{detect}, \code{fpr},
\code{published}), the planted-edge arm and its non-short-circuiting control, the \code{symnull}
and \code{abnull} true nulls, the \code{MEASURED\_SD} constants recorded 2026-08-31 at 1{,}200
pairs per condition, the \code{PUBLISHED} claim list, and the \code{NOMINAL ONLY --- underpowered}
verdict rule.

\bibitem{crossplay}
FishAI project.
\emph{CROSSPLAY.md --- FishAI against the fishlabs bots, in the fishlabs engine}.
Repository document; \S2 (the per-opponent cells at 150 games each, published without intervals),
\S3 (defusal replicated in a foreign engine, 300 duplicate pairs per opponent, half-widths
$0.517$--$0.560$).

\bibitem{bounded}
FishAI project.
\emph{BOUNDED.md --- the bit-budget memory ladder}.
Repository document; \S5a (the v1.5-against-v1.0 ladder, 2{,}000 duplicate pairs per rung on bank
\code{v15v10-holdout-a}, and the residual at an unbounded budget).

\bibitem{botlab}
FishAI project.
\emph{BOT\_LAB.md --- play-style spectrum: research, experimental design, and metrics}.
Repository document; \S5.1--5.2 (duplicate-deal pairing), \S5.4 (fixed-$N$ precision for published
estimates), \S5.5 (SPRT for tuning decisions only).

\bibitem{inertaxis}
A.~Hsieh.
\emph{The inert axis: when a style parameter is wired, swept, and never reached}.
Companion paper, FishAI project, 2026. The decision-level divergence certificate, and the doctrine
that a correction belongs adjacent to the claim it corrects.

\bibitem{containedbook}
A.~Hsieh.
\emph{The contained book: an absorbing resource measured to be worth nothing}.
Companion paper, FishAI project, 2026. The project's other principal null, reported in mean-score
units rather than in paired set-differences.

\bibitem{v05}
A.~Hsieh.
\emph{FishAI v0.5: measuring play style without skill in a 54-card Literature variant}.
Companion paper, FishAI project, 2026. The engine, the \usff{} dialect, the nine-style roster, and
the tournament design every instrument in this paper is built on.

\bibitem{v15}
A.~Hsieh.
\emph{FishAI v1.5: a bit-budget memory ladder as an honest difficulty axis}.
Companion paper, FishAI project, 2026. The generation the mostly-agreeing-arms discussion of
Section~\ref{sec:counter} is set against. It does \emph{not} supply that section's
\emph{v1.0}-against-\emph{v0.5} cell, which appears in no committed document
(Section~\ref{sec:untraceable}); the traceable mostly-agreeing-arms measurement on this ladder is
\cite{bounded}~\S5a's v1.5-against-v1.0 rung.

\bibitem{altman}
D.~G.~Altman and J.~M.~Bland.
Absence of evidence is not evidence of absence.
\emph{BMJ}, 311(7003):485, 1995.

\bibitem{cohen}
J.~Cohen.
\emph{Statistical Power Analysis for the Behavioral Sciences}, 2nd edition.
Lawrence Erlbaum Associates, Hillsdale, NJ, 1988.

\bibitem{gelmancarlin}
A.~Gelman and J.~Carlin.
Beyond power calculations: assessing Type S (sign) and Type M (magnitude) errors.
\emph{Perspectives on Psychological Science}, 9(6):641--651, 2014.

\bibitem{button}
K.~S.~Button, J.~P.~A.~Ioannidis, C.~Mokrysz, B.~A.~Nosek, J.~Flint, E.~S.~J.~Robinson, and
M.~R.~Munaf\`o.
Power failure: why small sample size undermines the reliability of neuroscience.
\emph{Nature Reviews Neuroscience}, 14(5):365--376, 2013.

\bibitem{ioannidis}
J.~P.~A.~Ioannidis.
Why most published research findings are false.
\emph{PLoS Medicine}, 2(8):e124, 2005.

\bibitem{amrhein}
V.~Amrhein, S.~Greenland, and B.~McShane.
Scientists rise up against statistical significance.
\emph{Nature}, 567:305--307, 2019.

\end{thebibliography}

\end{document}
